\section{Experimental Setup}
\label{sec:setup}

\subsection{Benchmarks}

\emph{In one sentence:} six benchmarks spanning knowledge, commonsense, mathematics, and logical reasoning supply both the degenerate and the high-entropy regimes the analysis requires.

TruthfulQA-MC1 \citep{lin2022truthfulqa} contributes adversarial factual questions with four or five options; we shuffle option order per question to avoid position bias. PIQA \citep{bisk2020piqa} contributes two-choice physical commonsense problems. LogiQA \citep{liu2020logiqa} contributes four-choice logical reading comprehension. LSAT logical reasoning and SAT mathematics, both from the AGIEval suite \citep{zhong2023agieval}, contribute five-choice analytical reasoning and four-choice quantitative items respectively. GSM8K \citep{cobbe2021training} contributes free-form grade-school mathematics problems answered with chain-of-thought reasoning and a final numeric value. We sample fifteen questions per benchmark uniformly at random with a fixed seed, giving ninety questions in total. For the held-out suite, a second draw of 119 fresh questions (twenty per benchmark, nineteen for LSAT logical reasoning) is taken from the same raw pools with the same filters by replaying the seeded shuffle and skipping consumed items, so the two sets share a construction procedure and nothing else. The mixture is deliberately diverse in difficulty: the strongest models in our pool saturate five of the six benchmarks while all models find LogiQA and, for some models, GSM8K and SAT mathematics challenging, which supplies both the degenerate and the high-entropy regimes that our analysis requires.

\subsection{Models and suites}

\emph{In one sentence:} ten commercial systems from five families, accessed through third-party free-tier endpoints in two phases --- a development suite on six systems and a held-out suite carried by four fully covered models.

The study evaluates ten commercial systems from five families, accessed through three providers under their free tiers and organized into the two question suites of Section~\ref{sec:method}. The \emph{development suite} (90 questions, collected October 17--18, 2026) covers six systems: Qwen3.8-Max and Qwen3.8-Omni from the Qwen family \citep{yang2025qwen3} on the Aliyun MaaS endpoint; Agnes3.0-Flash and Agnes2.5-Flash from the Agnes family \citep{agnesai2026} on the Agnes API endpoint; GPT-6-Luna from the GPT family \citep{openai2026gpt6} and Gemini3.8-Flash from the Gemini family \citep{google2026gemini38}, both on the Apinex endpoint. Its per-model coverage was truncated unevenly by free-tier quotas, to between 12 and 90 complete question matrices. The \emph{held-out suite} (119 fresh questions, collected October 19--23, 2026) re-evaluates the two Agnes models on fresh items and adds four DeepSeek-family systems \citep{deepseek2026api} and Qwen3.7-Flash: DeepSeekV4.1-Flash, first attempted on Aliyun MaaS (9 questions, retained as a provider record) and completed on the Apinex endpoint; DeepSeekV4-Pro-0813, a pro-tier sibling exposed by the same endpoint; DeepSeekV4-Flash-0731, stopped at 13 questions when its endpoint moved it behind a paid subscription; and Qwen3.7-Flash, stopped at 60 questions when its provider revoked free-tier access. Four systems complete all 119 questions and carry the primary held-out analysis; the two partial systems are reported separately and excluded from tabulation. The models span capability tiers from fast to flagship. All systems are queried through OpenAI-compatible endpoints in direct-answer mode on multiple-choice tasks and chain-of-thought on GSM8K; for the two Qwen flagships the proprietary thinking mode is disabled on multiple-choice tasks and enabled on GSM8K. Sampling uses temperature $0.7$ with default nucleus probability; greedy calls use temperature $0$. Exact endpoint model identifiers, per-call timestamps, and token usage are recorded in the logs released with the submission. Because the models are served through third-party endpoints --- Aliyun MaaS, the Agnes API, and the Apinex relay, the last an aggregator that served GPT-6-Luna, Gemini3.8-Flash, and the completed DeepSeek runs --- we can neither verify that a relay serves the requested model identity and version against the provider's first-party API nor confirm that the requested temperature is honored, especially for reasoning-style models; the Qwen thinking mode is the one decoding parameter we control explicitly, no condition in this study was replicated on a first-party API or an open-weight model, and findings therefore describe the accessed endpoints as of October 2026. We cite each provider's published technical report or model card where one exists and release the exact endpoint model identifiers, per-call timestamps, token usage, and raw outputs, so that the accessed configuration is auditable even where it cannot be verified\CHECK{verify every model citation below against the provider's official documentation for the exact served version; see the citation checklist in the submission notes}. Table~\ref{tab:models} states families, endpoints, access dates, parse rates, and completed coverage per suite.

\begin{table}[htbp]
\centering
\small
\caption{Evaluated systems, endpoints, and coverage across the two suites. Parse rate is the fraction of calls from which the mandated answer format was extracted after retries; $n$ is the number of questions with a complete call matrix. $^\dagger$Also the stem-paraphrase generator; its held-out row carries a self-preference risk (Section~\ref{sec:setup}) and is a partial replication excluded from pooled analyses. $^\ddagger$Completed on the second endpoint after the first exhausted its free tier. $^\S$Partial (13 of 119), excluded from tabulation. All systems run direct-answer on multiple-choice tasks and chain-of-thought on GSM8K.}
\label{tab:models}
\begin{tabular*}{\columnwidth}{@{\extracolsep{\fill}} l l l c c}
\toprule
\multicolumn{5}{l}{\textbf{Panel A: Development suite} (Oct 17--18, 2026)} \\
\midrule
Model & Family / endpoint & Tier & Parse rate & $n$ \\
\midrule
Qwen3.8-Max & Qwen / Aliyun & flagship & 99.6\% & 90 \\
Qwen3.8-Omni & Qwen / Aliyun & multimodal & 96.8\% & 71 \\
Agnes3.0-Flash & Agnes / Agnes & fast & 94.7\% & 33 \\
Agnes2.5-Flash & Agnes / Agnes & reasoning & 96.9\% & 40 \\
GPT-6-Luna & GPT / Apinex & reasoning & 100.0\% & 15 \\
Gemini3.8-Flash & Gemini / Apinex & fast & 99.6\% & 12 \\
\midrule
\multicolumn{5}{l}{\textbf{Panel B: Held-out suite} (Oct 19--23, 2026)} \\
\midrule
Model & Family / endpoint & Tier & Parse rate & $n$ \\
\midrule
Agnes3.0-Flash & Agnes / Agnes & fast & 98.7\% & 119 \\
Agnes2.5-Flash & Agnes / Agnes & reasoning & 98.1\% & 119 \\
Qwen3.7-Flash$^{\dagger}$ & Qwen / Aliyun & fast & 95.4\% & 60 \\
DeepSeekV4.1-Flash$^{\ddagger}$ & DeepSeek / Aliyun $\to$ Apinex & fast & 99.9\% & 119 \\
DeepSeekV4-Pro-0813 & DeepSeek / Apinex & pro & 99.9\% & 119 \\
DeepSeekV4-Flash-0731$^{\S}$ & DeepSeek / Aliyun & fast & 93.1\% & 13 \\
\bottomrule
\end{tabular*}

\end{table}

\subsection{Prompt design and paraphrase generation}

\emph{In one sentence:} paraphrases target the parts of the prompt where semantic-preserving variation is well defined --- stems on three datasets, instructions on the other three --- and pass a three-stage filter plus a bidirectional entailment check.

The canonical prompt for a multiple-choice question consists of the passage when present, the question stem, the option list labeled $(A)$ through $(E)$, and a fixed instruction that requires a direct answer in the form \texttt{Answer: <letter>} without explanation. GSM8K prompts instead require step-by-step reasoning ending with \texttt{Answer: <number>}. Direct answering on multiple-choice tasks keeps the comparison uniform across models, keeps token costs bounded, and isolates the aggregation question from chain-of-thought length effects; the GSM8K condition covers the reasoning-trace regime.

Paraphrases target the parts of the prompt where semantic-preserving variation is well defined: the question stem for TruthfulQA, LSAT, and LogiQA, and the instruction wording for GSM8K, PIQA, and SAT mathematics, whose stems contain problem-specific content such as numbers that should not be altered. This design yields two paraphrase loci, and all fragility and gain statistics are reported stratified by locus. Stem paraphrases are generated by Qwen3.7-Flash, prompted to rewrite each stem in six ways while preserving all facts, entities, numbers, and the question type; in the development suite this generator is not part of the evaluation pool. The held-out suite later evaluates Qwen3.7-Flash as a target system, a dual role we flag here and wherever its rows appear: its paraphrase-condition results carry a self-preference risk, since the model judges its own rewrites, and its row is moreover a partial replication (60 of 119 questions) excluded from the pooled analyses. Instruction paraphrases are six handcrafted variants per dataset that vary wording while keeping the literal answer-format requirement fixed, which keeps answer extraction reliable. As a robustness check, a second generator (Agnes3.0-Flash, outside its own evaluation role) regenerates paraphrases for all stem-locus questions, and the alternative set is validated to the same standard; a planned re-evaluation of the prompt-brittle models under the alternative paraphrases was cut short by provider quota exhaustion and is left to future work.

Every generated paraphrase passes a three-stage validation filter. A lexical stage rejects paraphrases that are identical to the original or whose length ratio falls outside $[0.4, 2.5]$. A semantic stage embeds the original stem and each candidate with a text embedding model and rejects candidates whose cosine similarity to the original falls below $0.75$. A form stage, applied to factual questions, rejects rewrites that change the question type. Candidates that fail any stage are replaced by regenerated candidates until six valid paraphrases per question are obtained; the final set achieves a mean embedding similarity of $0.86$ with the originals. One further check strengthens the semantic stage: a bidirectional entailment test, administered by an auxiliary language model, asks whether each paraphrase entails the original and is entailed by it; 267 of 270 final paraphrases pass both directions, and 269 of 270 of the second generator's paraphrases pass, with the full per-item verdicts released alongside the paraphrases. As a sensitivity analysis, the fragility statistics are recomputed under a stricter embedding threshold of $0.80$, and the qualitative pattern is unchanged. These automatic filters bound equivalence only in distribution; they do not certify item-level semantic equivalence, a limitation we state in Section~\ref{sec:discussion} and address with a planned human spot-check (Section~\ref{sec:discussion}).

\subsection{Execution and quality control}

\emph{In one sentence:} an asynchronous runner with rate limiting, retries, and append-only checkpointing collects the sixteen-call matrix; questions with any unparsed cell are excluded from that model's analysis, and two sensitivity analyses bound the effect of the exclusions.

The experiment executes the sixteen-call matrix of Section~\ref{sec:method} for every question and model through an asynchronous runner with per-provider rate limiting, automatic retry with exponential backoff, and append-only checkpointing. Answer extraction uses the mandated \texttt{Answer:} format first and falls back to scanning the final output lines for a well-formed option letter or number, with the option range constrained by the question. Every stored record includes the extracted answer, its correctness, token usage, and a content snippet for auditing. Extraction succeeds on between 94.7\% and 100\% of calls across the ten systems (Table~\ref{tab:models}); cells whose first decode fails to parse are re-drawn under the identical prompt and decoding configuration until extraction succeeds or the free-tier quota of the provider is exhausted, and questions with any remaining unparsed cell are excluded from that model's analysis, so all reported accuracies are computed over complete question matrices. Two safeguards address the residual exclusions: a lenient-aggregation sensitivity analysis that admits any question with at least four parsed votes per source, and a comparison of excluded with included questions on greedy accuracy to probe for difficulty-related selection bias; both are reported in \ref{app:dev}. Per-model completed question counts, which vary between 12 and 90 on the development suite, are stated alongside every result, and pooled analyses use the intersection of completed questions across the models involved, with question-level cluster bootstrap intervals.

\subsection{Follow-up conditions}

\emph{In one sentence:} three follow-up conditions --- greedy paraphrases, an MMLU-Pro probe, and chain-of-thought answering --- reuse the main matrix's prompts and parsing and are analyzed in the appendices as descriptive robustness checks.

\textbf{Greedy paraphrases.} Each of the six paraphrases is decoded once at $\tau = 0$, yielding the noise-free paraphrase ensemble $\mathrm{PE}^{g}\text{-}6$ (\ref{app:greedypara}); the control covers five of the six development-suite models and was \emph{not} re-collected on the held-out suite, whose remaining provider budget went to the fresh-question matrices. Held-out paraphrase conditions therefore retain the decoding-noise mixing that this control removes, and are read with that caveat (\ref{app:greedypara}).

\textbf{Hard-task probe.} Thirty questions sampled uniformly at random (fixed seed, stratified by subject) from the MMLU-Pro test set \citep{wang2024mmlupro}, administered with the greedy baseline and the six-sample self-consistency condition under direct answering (\ref{app:mmlupro}).

\textbf{Chain-of-thought answering.} The greedy baseline, the six-sample self-consistency condition, and, for stem-locus questions, the six-paraphrase ensemble, repeated under an explicit step-by-step instruction (\ref{app:cot}); GSM8K is excluded because its main-suite condition is already chain-of-thought.

\subsection{Conditions and metrics}

\emph{In one sentence:} each model and question yields the baselines, ensembles, agreement statistics, modal share, and fragility indicator defined in Section~\ref{sec:method}; accuracy is the fraction of questions answered correctly, error detection is evaluated by AUROC and tie-randomized AURC, and budget curves plot accuracy against $k$. All analysis code operates on the raw prediction logs, which are released with the submission together with the prompts and paraphrases.
